\subsection{Iterating the deleted-line
residue}\label{iterating-the-deleted-line-residue}

Write \(q=p^e>3\), and let \(L\) be the line \(y=0\) in \(AG(2,q)\). The
source's rank-two residue \[
S_q=\Gamma(AG(2,q))_L
\] has the following literal description. Its points are \((x,y)\) with
\(y\ne0\). Its lines are \[
M_{a,b}:\ x=ay+b\qquad(a,b\in\mathbb F_q),
\] with their points on \(L\) removed. Each line has \(q-1\) points, and
each point belongs to \(q\) lines. This is exactly the residue of
Construction 3: a flag \((z,M)\) incident to \(L\) has
\(M\cap L=\{z\}\), so its line \(M\) is nonparallel to \(L\) and its
intersection \(z\) is uniquely determined.

Pass to the dual \(N_q=S_q^*\). Its points are the pairs
\((a,b)\in\mathbb F_q^2\). For an original point \((x,y)\) with
\(y\ne0\), its dual line is \[
b=-ya+x.
\] Thus \(N_q\) is the affine plane with both the horizontal and the
vertical parallel classes omitted. Its lines are exactly \[
\ell_{m,c}:\ b=ma+c\qquad(m\in\mathbb F_q^*,\ c\in\mathbb F_q).
\tag{D1}
\] There are \(q^2\) points and \(q(q-1)\) lines. The line size is
\(q\), and the number of lines on a point is \(q-1\).

We use Construction3 literally. For any partial linear space \(T\), its
objects are the points \(P\), lines \(\ell\), and incident pairs
\(F=(x,M)\) of \(T\), with \[
P\sim\ell\iff P\notin\ell,\quad
P\sim(x,M)\iff P\in M\setminus\{x\},\quad
\ell\sim(x,M)\iff\ell\cap M=\{x\}.
\tag{D2}
\] The construction commutes with duality, exchanging its first two
types. To verify the duality assertion of the source's Lemma 2, send an
old point to its dual line, an old line to its dual point, and \((x,M)\)
to \((M^*,x^*)\). The first relation in (D2) is symmetric under this
exchange. The second becomes the assertion that two dual lines meet in
the indicated dual point, which is equivalent to uniqueness of the
original line through two distinct points. The third becomes membership
of the indicated original point in a distinct original line. These
equivalences use precisely the partial-linearity assumption. Hence \[
\Gamma(S_q)\cong\Gamma(N_q).
\tag{D3}
\]

\subsubsection{Its full automorphism
group}\label{its-full-automorphism-group}

All automorphism groups in this note preserve the displayed types.

\textbf{Proposition D.Aut.} Every automorphism of \(N_q\) is of one of
the forms \[
(a,b)\longmapsto
(\alpha a^\sigma+u,\ \beta b^\sigma+v)
\quad\hbox{or}\quad
(\alpha b^\sigma+u,\ \beta a^\sigma+v),
\tag{D4}
\] where \(\alpha,\beta\in\mathbb F_q^*\), \(u,v\in\mathbb F_q\), and
\(\sigma\in\operatorname{Gal}(\mathbb F_q/\mathbb F_p)\). In particular
\[
|\operatorname{Aut}(S_q)|=|\operatorname{Aut}(N_q)|
=2e\,q^2(q-1)^2.
\tag{D5}
\] Restriction also gives \[
\operatorname{Aut}(\Gamma(N_q))\cong\operatorname{Aut}(N_q).
\tag{D6}
\]

\textbf{Proof.} Two distinct points of \(N_q\) fail to be collinear
exactly when they have the same first coordinate or the same second
coordinate. The noncollinearity graph is therefore the Hamming graph on
\(\mathbb F_q^2\). Its maximal cliques are exactly the rows and columns:
once two vertices of a clique have the same first coordinate and
different second coordinates, a vertex adjacent to both must share that
first coordinate. The assertion with the coordinates reversed is
identical.

The rows form one family of \(q\) pairwise disjoint maximal cliques, and
the columns the other; every row meets every column in one point.
Consequently an automorphism either permutes each family or exchanges
the two. After an optional coordinate exchange, its point permutation
has the form \[
(a,b)\longmapsto(\phi(a),\psi(b))
\] for arbitrary permutations \(\phi,\psi\) of the field.

Translations and independent nonzero coordinate scalings are actual
automorphisms of (D1). Compose with such maps to arrange \[
\phi(0)=\psi(0)=0,\qquad \phi(1)=\psi(1)=1.
\] The diagonal is then fixed, being the unique line through \((0,0)\)
and \((1,1)\); hence \(\phi=\psi=\theta\).

Parallelism of lines is intrinsic as disjointness. Since the diagonal is
fixed, its entire parallel class is preserved. The line \(b=a+c\) must
therefore map to \(b=a+\theta(c)\), its intercept being the image of
\((0,c)\). It follows that \[
\theta(a+c)=\theta(a)+\theta(c).
\] Likewise the line \(b=ma\) maps to the line through \((0,0)\) and
\((1,\theta(m))\), giving \[
\theta(ma)=\theta(m)\theta(a).
\] These formulas also hold at zero. Thus \(\theta\) is a field
automorphism. This proves necessity in (D4); substitution in (D1) proves
sufficiency. The parameters are unique, and the exchanged and
unexchanged forms are distinct for \(q>2\). Counting them gives (D5).

Finally any type-preserving automorphism of \(\Gamma(N_q)\) restricts on
its first two types to an automorphism of \(N_q\), since original
incidence is the complement of their incidence in (D2). The
point-neighbors of a third-type object \((x,M)\) form
\(M\setminus\{x\}\). This set contains at least two points, which
uniquely recover \(M\), and its complement in \(M\) then recovers \(x\).
Thus restriction is injective. Every automorphism of \(N_q\) plainly
extends to its flags while preserving (D2), proving surjectivity and
(D6). ∎

The source's preceding assertion on p.568 identifies
\(\operatorname{Aut}(S_q)\) with the affine-plane stabilizer of \(L\),
whose order is \(e q^2(q-1)^2\). Proposition D.Aut exhibits the missing
factor two. In the dual model it is the exchange of coordinates. This
exchange already exists before Construction3 is applied a second time.

\subsubsection{Firmness and rank-two
incidence}\label{firmness-and-rank-two-incidence}

\textbf{Proposition D.Firm.} For every \(q>3\), \(\Gamma(N_q)\) is a
firm geometry and satisfies \((IP)_2\). It has \[
q^2,\quad q(q-1),\quad q^2(q-1)
\] objects of its three respective types, and \[
q^2(q-1)^2(q-2)
\tag{D7}
\] chambers.

\textbf{Proof.} Every incident pair in (D2) extends to a chamber. For
\(P\notin\ell\), choose a line \(M\) through \(P\) with slope different
from that of \(\ell\). There are exactly \(q-2\) such choices; each
intersects \(\ell\) at a unique point \(x\ne P\) and supplies the flag
\((x,M)\). For \(P\in M\setminus\{x\}\), choose any of the \(q-2\) lines
through \(x\) other than \(M\). For \(\ell\cap M=\{x\}\), choose any of
the \(q-1\) points of \(M\setminus\{x\}\). Thus the three types of
incident pair lie in respectively \(q-2\), \(q-2\) and \(q-1\) chambers.
All these numbers are at least two. Single objects and the empty flag
extend because each object has incident pairs. This proves both the
geometry axiom and firmness.

A chamber is uniquely specified by \((x,M)\), a point
\(P\in M\setminus\{x\}\), and a line \(\ell\ne M\) through \(x\). The
three choices have sizes \(q^2(q-1)\), \(q-1\) and \(q-2\), proving
(D7).

The residue of a flag object \((x,M)\) is the complete bipartite graph
on the \(q-1\) points of \(M\setminus\{x\}\) and the \(q-2\) other lines
through \(x\). For a point residue, its flag objects correspond to the
distinct points collinear with that point, and a pair of them lies on at
most one line of \(N_q\). Thus this residue is a partial linear space.
The analogous statement for a line residue follows by duality, or
directly from the fact that two distinct lines meet in at most one
point. Hence every rank-two residue is either a partial linear space or
a generalized digon, exactly \((IP)_2\). ∎

\subsubsection{Residual connectedness}\label{residual-connectedness}

For a partial linear space \(T\) and point \(z\), form the graph on the
points collinear with \(z\), with an edge between \(x,y\) when \(x,y,z\)
are pairwise collinear but are not on one line. In the point residue
\(\Gamma(T)_z\), its flag objects are naturally these collinear points:
\((x,M)\) is uniquely determined by \(x\) because \(M\) is the unique
line through \(x,z\). Its other objects are the lines not on \(z\). Such
a line is incident to all its points collinear with \(z\). Consequently
the displayed graph is the graph obtained by replacing each of these
line objects by a clique on its neighbors. Provided no such line object
is isolated, connectivity of the two graphs is equivalent. We apply this
elementary reduction to \(N_q\) and its dual, paying the non-isolation
condition explicitly below.

\textbf{Proposition D.Connected.} For \(q>3\), \(\Gamma(N_q)\) is
residually connected if and only if \(q>4\). At \(q=4\), each residue of
an original point of \(N_q\) is the disjoint union of three hexagons.

\textbf{Proof.} First take the original point \(z=(0,0)\) of \(N_q\).
Its collinear points are \[
(a,b)\in(\mathbb F_q^*)^2.
\] Two such points are adjacent in the preceding graph exactly when \[
a\ne a',\qquad b\ne b',\qquad b/a\ne b'/a'.
\tag{D8}
\] The first two inequalities give their collinearity in the net, and
the last says that their joining line does not contain the origin. Every
net line avoiding the origin has exactly \(q-2\) points collinear with
the origin: its intersections with the horizontal and vertical axes are
distinct, and all its other points have both coordinates nonzero. Thus
no line object of the residue is isolated.

Choose a generator \(\zeta\) of \(\mathbb F_q^*\) and put \(n=q-1\). In
exponents, (D8) is the Cayley graph of \[
(\mathbb Z/n\mathbb Z)^2
\] whose allowed differences \((r,s)\) satisfy \(r\ne0\), \(s\ne0\) and
\(r\ne s\). If \(n\ge4\), both \((1,2)\) and \((1,3)\) are allowed.
Their difference is \((0,1)\) and subtracting twice this vector from
\((1,2)\) gives \((1,0)\). They therefore generate the whole group, so
the graph and the original point residue are connected.

For \(q=4\), one has \(n=3\). The only allowed differences are \((1,2)\)
and \((2,1)\), which generate the subgroup of order three on which the
sum of coordinates is zero. The graph has three components, each a
triangle. Every line object of the incidence residue has degree two, and
each pair of adjacent point objects has a unique joining line. Each
triangle therefore lifts to a hexagon. There are exactly three disjoint
hexagons, not one connected residue.

Next take an original line of \(N_q\). By the coordinate automorphisms
it suffices to use the diagonal \(\ell:b=a\). In the dual near-point
graph its vertices are the lines of \(N_q\) meeting \(\ell\), other than
\(\ell\) itself. Write such a line by its slope
\(m\in\mathbb F_q^*\setminus\{1\}\) and its intersection \((t,t)\) with
\(\ell\). Two vertices \((t,m),(t',m')\) form a proper dual triangle
with \(\ell\) exactly when \[
t\ne t'\quad\hbox{and}\quad m\ne m'.
\tag{D9}
\] If the slopes differ their intersection exists; it lies on \(\ell\)
precisely when the two displayed intersection points agree. Thus this
graph is the direct product \(K_q\times K_{q-2}\). It is connected for
\(q>3\). One elementary verification is that two vertices agreeing only
in the second coordinate have a two-edge path through a new first
coordinate and a different second coordinate. For vertices agreeing only
in the first coordinate the same argument applies if \(q-2\ge3\); when
\(q-2=2\), a three-edge path using two other first coordinates suffices.

The original point objects in this line residue correspond to points
outside \(\ell\). Through each such point there are \(q-2\) net lines
meeting \(\ell\): exclude only the parallel direction. Hence the clique
expansion has no isolated omitted objects, and the line residue is
connected.

The remaining rank-two residues are the complete bipartite flag residues
from Proposition D.Firm, so they are connected. The full incidence graph
is connected as well. Indeed any two original points have an original
line avoiding both, since at most \(2(q-1)\) of the \(q(q-1)\) lines
contain either point. Any two original lines have an original point
outside their union, since \(q^2>2q\). Thus the subgraph on the first
two types is connected, and every flag object has neighbors there. These
are all residues of rank at least two in a rank-three geometry,
completing the proof. ∎

Under (D3), the disconnected residues at \(q=4\) are exactly the
original-line residues of \(\Gamma(S_4)\). The source's claimed
exceptional geometry has 12,16,48 objects and 288 chambers, but its
sixteen original-line residues each have three hexagonal components. Its
failure is residual connectedness, independently of any group-order
calculation.

\subsubsection{Flag transitivity}\label{flag-transitivity}

\textbf{Proposition D.Transitivity.} For \(q=p^e>3\), \(\Gamma(N_q)\) is
flag-transitive if and only if \(q=4\) or \(q=8\).

\textbf{Proof.} A chamber can be encoded as the ordered configuration \[
P\in M,\quad x\in M,\quad P\ne x,\quad x\in\ell,\quad M\ne\ell.
\] The group (D4) can translate \(x\) to \((0,0)\) and independently
scale the two coordinates to send \(P\) to \((1,1)\). Both coordinate
differences are nonzero because \(M\) is a net line. This also sends
\(M\) to the diagonal. The remaining line \(\ell\) then has an arbitrary
slope \[
\lambda\in\mathbb F_q^*\setminus\{1\}.
\] The stabilizer of this normalized ordered data consists exactly of
simultaneous field automorphisms, optionally followed by coordinate
exchange. It acts on \(\lambda\) by \[
\lambda\longmapsto\lambda^{p^j},\qquad
\lambda\longmapsto\lambda^{-p^j}.
\tag{D10}
\] This follows directly from (D4): fixing the origin removes both
translations, and fixing \((1,1)\) sets both scaling constants equal to
one. Hence chamber transitivity is equivalent to transitivity of the
action (D10).

Its acting group has order dividing \(2e\). Therefore a necessary
condition is \[
q-2\mid2e.
\tag{D11}
\] For prime powers \(q>3\), this leaves only \(q=4,8\). For \(p\ge5\),
already \(p^e-2>2e\) for all \(e\ge1\). For \(p=3\), the range \(q>3\)
gives \(e\ge2\), and \(3^e-2>2e\). For \(p=2\), all \(e\ge4\) satisfy
\(2^e-2>2e\). These inequalities follow from their first case by
induction. The two remaining powers of two are four and eight.

For \(q=4\), inversion exchanges the two elements of
\(\mathbb F_4^*\setminus\{1\}\). For \(q=8\), choose a generator
\(\zeta\) of the multiplicative group of order seven. The six exponents
\[
1,\ 2,\ 4,\ -1,\ -2,\ -4\pmod7
\] are all the nonzero residues modulo seven. Thus the Frobenius and
inversion orbit of \(\zeta\) is exactly \(\mathbb F_8^*\setminus\{1\}\).
This proves sufficiency in both cases. By Proposition D.Aut these are
the full type-preserving automorphisms of the constructed geometry, so
no additional chamber transitivity was omitted. ∎

\subsubsection{The corrected exceptional
value}\label{the-corrected-exceptional-value}

\textbf{Theorem D.DeletedLine.} Let \(q>3\) be a prime power, let \(L\)
be a line of \(AG(2,q)\), and put \(S_q=\Gamma(AG(2,q))_L\) using the
source's Construction3. Then:

\begin{enumerate}
\def\labelenumi{\arabic{enumi}.}
\tightlist
\item
  \(\Gamma(S_q)\) is a firm geometry satisfying \((IP)_2\) for every
  such \(q\).
\item
  It is residually connected exactly when \(q>4\).
\item
  It is flag-transitive exactly when \(q=4\) or \(q=8\).
\item
  Consequently it is an A-geometry exactly when \(q=8\).
\end{enumerate}

At \(q=8\) the three type sizes, in the source's original
point/line/flag order, are 56,64,448. There are 18,816 chambers, and its
full type-preserving automorphism group has order 18,816 and acts
regularly on them. The group is the dual of the semiaffine group in
(D4).

\textbf{Proof.} Combine (D3) and the four preceding propositions. The
count and regularity assertions follow from (D5), (D7) and transitivity.
∎
